\chapter{Introduction}

\section{Dialectic Debating of Related Concepts —— Analysis of Contradictions in the Context of Computation}
\label{sec:dialectics}

To open this thesis, I want to begin not with a research question but with a debate. The object I am studying---a reconfigurable inflatable modular robot---sits at the intersection of several long-running tensions in how we think about machines, materials, and computation. Rather than picking a side too quickly, it is worth laying these tensions out as dialectical pairs: seemingly opposed, but in practice deeply entangled. The point is not to elaborate each pair exhaustively, but to extract the underlying axes along which my research will eventually have to take a position.

\sidebyside{Natural vs. Artificial}{
Everything originates from nature, yet human intelligence introduced a peculiar rupture: the capacity to \textit{produce artifacts}. This marks the emergence of the artificial---not as something outside nature, but as a recursive extension of it. Antoni Gaud\'i's claim---\textit{``straight lines belong to men, curves to God''}---reveals a deeper distinction: artificial systems historically favor abstraction, regularity, and control, while natural systems embody continuous variation, adaptation, and emergence. Yet this boundary is increasingly unstable. Bio-inspired design, generative fabrication, and soft robotics suggest that artificial systems no longer merely sit \textit{against} nature but can \textit{grow, adapt, and behave} like natural ones. The artificial, in this sense, becomes a \textit{second-order nature}---a designed ecology of behaviors and materials.
}

\sidebyside{Physical vs. Digital}{
Before the digital computer, computation was largely \textit{physical}. Analog machines---radios, phonographs, and Lord Kelvin's tide-predicting machine of 1872---encoded values not as symbols but as continuous changes in physical quantities such as voltage, pressure, or position. The pulleys and gears of the tide machine were not representing the equations; they \textit{were} the equations. The digital turn replaced this with discrete symbols, gaining repeatability and verifiability while installing a particular metaphor at the center of how we think about machines: that intelligent behavior must come from a controller running symbols, with the body merely as the thing the controller commands. Recent work on \textit{morphological computation} \cite{pfeifer2007body} pushes back, arguing that the body itself was already doing much of the computing the digital metaphor abstracted away.
}

\sidebyside{Predictable vs. Unpredictable}{
Digital signals are typically discrete and explicit---a 1 or a 0, on or off. Such normative digits are easy to manipulate, easy to verify, and have become the fundamental units of computation as we typically practice it. Predictability is comfortable: it allows proofs, tests, and guarantees. Physical systems, by contrast, tend toward the unpredictable---not in the sense of being random, but in the sense of being \textit{underdetermined} by any model we are willing to write down. A balloon does not inflate to exactly the same shape twice. Yet there is an older tradition---Turing's reaction--diffusion patterns, the foraging trails of \textit{Physarum polycephalum}, ant colony optimization---that treats this unpredictability not as noise to be rejected but as a \textit{resource} from which interesting behavior emerges.
}

\sidebyside{Bottom-up vs. Top-down}{
If predictability is the epistemic axis, bottom-up versus top-down is the design axis. A top-down system specifies the global behavior first and decomposes it into commands for each component; a bottom-up system specifies only \textit{local} rules and lets global behavior arise from their interaction. Cellular automata, flocking models, swarm intelligence, and the self-assembly of programmable matter \cite{rubenstein2014programmable} live on the bottom-up side. Top-down systems are tractable and easy to debug; bottom-up systems are robust, scalable, and degrade gracefully. Neither is universally better. The interesting cases---and, as I will argue, the case of an inflatable modular robot whose modules are physically coupled through their material---are the ones where neither approach alone is sufficient.
}

\sidebyside{Soft vs. Rigid}{
Soft and rigid is, in some sense, the physical instantiation of the dialectics above. Rigid bodies are predictable in shape, well-modeled by classical mechanics, and amenable to top-down control---which is why nearly all of industrial robotics is built on them. Soft bodies trade that predictability for compliance: they conform to objects they touch, absorb shocks, and continuously change their effective shape under load \cite{rus2015design}. Crucially, this compliance is not just a property to be tolerated; it can be a property to be \textit{recruited}. A soft gripper may grasp an irregular object using nothing but its own deformation, with no closed-loop feedback at all. The morphology, in this case, is doing the computation a controller would otherwise have to do. The soft/rigid distinction is therefore more than a material choice---it is a choice about \textit{where the intelligence lives}.
}

\section{Synthesis: From Dialectics to Design Space —— Where These Pairs Converge}
\label{sec:synthesis}

Laid side by side, these five pairs are not independent axes. They cluster. The \textit{natural}, the \textit{physical}, the \textit{unpredictable}, the \textit{bottom-up}, and the \textit{soft} all share a common posture: they treat behavior as something that \textit{arises} from a substrate rather than something \textit{imposed} on it. The \textit{artificial}, the \textit{digital}, the \textit{predictable}, the \textit{top-down}, and the \textit{rigid} share the opposite posture: behavior is specified, encoded, and executed. Most of contemporary robotics has been built on this second cluster, and within its domain it has been enormously successful. But a growing body of work---soft robotics \cite{rus2015design}, morphological computation \cite{pfeifer2005new}, modular self-reconfigurable systems \cite{yim2007modular}, programmable matter \cite{cademartiri2015programmable}, and particle robots \cite{li2019particle}---suggests that the next class of capabilities (adaptive, embodied, energetically efficient, robust to the messiness of the real world) requires us to take the first cluster seriously again, not as a return to pre-industrial craft but as a deliberate design strategy.

The problem is that these two clusters have not been brought together evenly. Soft robotics has produced compelling demonstrations of compliance and material intelligence, but most soft robots are \textit{monolithic}---a single continuous body whose behavior is hard to compare across designs. Modular and reconfigurable robotics has produced rich design spaces of comparable units, but those units are almost universally \textit{rigid}, and their reconfigurability lives in joint angles or magnetic re-attachment rather than in the body itself. Robots that are simultaneously \textit{soft, modular, and reconfigurable}---where the modules are compliant, where the body's shape is part of the computation, and where different topological arrangements can be compared on equal terms---remain a thin slice of the literature. This is the slice this thesis works in.

A reconfigurable inflatable modular robot sits squarely at the convergence of the dialectics above, which is why it is interesting as a design object. It is artificial, but its behavior is shaped by physical phenomena---air pressure, membrane elasticity, contact between neighboring inflated bodies---that resist neat symbolic description. It is built from discrete modules, but its motion is continuous and underdetermined. It is amenable to top-down command at the level of which valves to open, but its actual locomotion emerges from the bottom-up interaction of inflated neighbors and gravity. It is soft in its surface and rigid in its topology. Almost every dialectic above passes through this object.

Two consequences follow, and together they define the problem space of this thesis. First, \textit{configuration matters in a way it does not for rigid modular robots}. When the modules are soft and physically coupled, the geometric arrangement of modules is not just a mounting decision---it conditions how forces propagate, how neighbors deform one another, and ultimately what kinds of motion are possible. The configuration is part of the computation. Second, \textit{control cannot be inherited wholesale from either tradition}. Top-down gait planning struggles with a body that is too compliant to fully command; purely bottom-up local rules struggle to satisfy goal-directed tasks like ``move toward that target.'' Some negotiation between the two is required, and the form that negotiation takes depends on the configuration the robot has.

These two consequences together circumscribe a coherent design space: a family of robots built from identical inflatable modules, arranged in different but comparable geometric configurations, in which the relationship between \textit{shape}, \textit{actuation}, and \textit{controllability} can be studied as an interconnected whole rather than as separate engineering problems. The chapters that follow take this design space as their starting condition and ask, in different ways, what it means to design and to control a robot whose body is not a passive container for computation but an active participant in it.
